\documentclass[11pt]{article}
\usepackage[a4paper,margin=1in]{geometry}
\usepackage{fontspec}
\usepackage[boldfont=false]{xeCJK}
\setCJKmainfont{FandolSong-Regular.otf}
\usepackage{amsmath}
\usepackage{amssymb}
\usepackage{array}
\usepackage{booktabs}
\usepackage{graphicx}
\usepackage{adjustbox}
\usepackage{enumitem}
\setlist[itemize]{topsep=2pt,partopsep=0pt,parsep=0pt,itemsep=2pt,leftmargin=1.4em}
\setlist[enumerate]{topsep=2pt,partopsep=0pt,parsep=0pt,itemsep=2pt,leftmargin=1.6em}
\usepackage{parskip}
\usepackage{fvextra}
\fvset{breaklines=true,breakanywhere=true,fontsize=\small}
\setlength{\parindent}{0pt}

\begin{document}

\begin{center}
  {\LARGE\bfseries Chemical energetics}\\[0.35em]
  {\large A Level Chemistry}
\end{center}
\vspace{0.6em}

\section*{Enthalpy change, ΔH}

Every chemical reaction takes in or gives out energy. This energy change, measured at constant pressure, is the \textbf{enthalpy change} 焓变, with symbol $\Delta H$.

\begin{itemize}
\item in an \textbf{exothermic} 放热 reaction the system gives out heat, so the products have less energy than the reactants and $\Delta H$ is \textbf{negative}.

\item in an \textbf{endothermic} 吸热 reaction the system takes in heat, so the products have more energy than the reactants and $\Delta H$ is \textbf{positive}.

\end{itemize}

\subsection*{Reaction pathway diagrams}

A \textbf{reaction pathway diagram} 反应路径图 shows the energy of the reactants and products, and the energy "hill" between them. The height of the hill is the \textbf{activation energy} 活化能 — the least energy the particles need before they can react.

\begin{itemize}
\item exothermic: products sit \textbf{lower} than reactants ($\Delta H < 0$).

\item endothermic: products sit \textbf{higher} than reactants ($\Delta H > 0$).

\end{itemize}

\subsection*{Standard conditions and types of enthalpy change}

Energy values are compared under \textbf{standard conditions} 标准条件: $298\ \text{K}$ and $101\ \text{kPa}$, shown by the symbol $^{\ominus}$. Each substance is in its normal physical state at those conditions.

\begin{center}
\adjustbox{max width=\linewidth}{%
\begin{tabular}{lll}
\toprule
\textbf{Symbol} & \textbf{Name} & \textbf{Definition (per mole, under standard conditions)} \\
\midrule
$\Delta H_r^{\ominus}$ & \textbf{enthalpy change of reaction} 反应焓变 & for the amounts shown in the equation \\
$\Delta H_f^{\ominus}$ & \textbf{enthalpy change of formation} 生成焓变 & one mole of a compound forms from its elements \\
$\Delta H_c^{\ominus}$ & \textbf{enthalpy change of combustion} 燃烧焓变 & one mole of a substance burns completely in oxygen \\
$\Delta H_{\text{neut}}^{\ominus}$ & \textbf{enthalpy change of neutralisation} 中和焓变 & one mole of water forms from an acid and an alkali \\
\bottomrule
\end{tabular}}
\end{center}

\subsection*{Energy from breaking and making bonds}

During a reaction, old bonds break and new bonds form. \textbf{Breaking} a bond needs energy (endothermic); \textbf{making} a bond releases energy (exothermic). The enthalpy change of the reaction is the difference between the two:

\[
\Delta H_r = \sum (\text{bond energies broken}) - \sum (\text{bond energies made})
\]

The \textbf{bond energy} 键能 is the energy needed to break one mole of a particular bond in the gas state, so it is always positive. Some bond energies are exact (for one specific molecule); others are averages taken over many different molecules, so calculations using them are only approximate.

\subsection*{Measuring enthalpy change in the lab}

When a reaction heats up (or cools down) a known mass of water or solution, the heat transferred is:

\[
q = mc\Delta T
\]

where $m$ is the mass, $c$ is the \textbf{specific heat capacity} 比热容 (how much energy raises 1 g by 1 K), and $\Delta T$ is the temperature change. The enthalpy change per mole is then:

\[
\Delta H = -\frac{mc\Delta T}{n}
\]

The minus sign makes $\Delta H$ negative when the temperature rises (an exothermic reaction).

\section*{Hess's law}

\textbf{Hess's law} 盖斯定律 says that the total enthalpy change for a reaction is the same, no matter which route you take from reactants to products. This is because energy is conserved.

This lets you draw an \textbf{energy cycle} 能量循环: link the reactants and products by a direct step and by an indirect route, then add the steps so that both routes give the same total.

Hess's law is useful in two ways:

\begin{itemize}
\item it lets you find an enthalpy change that you \textbf{cannot} measure directly (for example, the formation of a compound that forms slowly or with side reactions).

\item it lets you calculate $\Delta H_r$ from bond energy data, or from formation or combustion data given in the question.

\end{itemize}


\end{document}
